% 01-methods.tex — Methods
% Software versions and parameters: data/Qiime2_Pipeline.sh
% Filtering and DADA2 counts: data/tables/00_filtering_funnel.tsv, 01_dada2_stats.tsv, 03_lme_model_summary.tsv

\section{Methods}

\subsection{Study design, data source, and cohort filtering}
This study reanalyzed publicly available longitudinal 16S rRNA gene amplicon sequencing data from a randomized clinical trial investigating autologous fecal microbiota transplantation (auto-FMT) in patients undergoing allogeneic hematopoietic cell transplantation (allo-HCT; ClinicalTrials.gov identifier NCT02269150) \citep{taur2018reconstitution}, distributed through the QIIME 2 Cancer Microbiome Intervention Tutorial. All bioinformatics and statistical workflows were executed in QIIME 2 (version 2026.7) \citep{bolyen2019reproducible} using a two-stage educational and analytical pipeline (\texttt{data/Qiime2\_Pipeline.sh}). Because the full study repository contained 12,546 fecal samples comprising 17,865 amplicon sequence variants (ASVs) and 550,243,459 total reads, an upstream module first demonstrated raw paired-end read import and denoising on a standalone subset of 41 demultiplexed fastq samples, whereas the main downstream analysis operated directly on the full pre-computed feature table (\texttt{feature-table.qza}) and representative sequences (\texttt{rep-seqs.qza}). Note that the upstream 41-sample demonstration and the downstream 12,546-sample cohort analysis are independent workflow segments and do not feed into one another.

To isolate the randomized auto-FMT trial cohort and focus on the acute peri-transplant recovery window, the downstream feature table underwent a five-step sequential filtering funnel (\texttt{qiime feature-table filter-samples}, \texttt{filter-features}, and \texttt{qiime taxa filter-table}):
\begin{enumerate}
  \item \textbf{Trial cohort selection:} Samples were restricted to participants assigned to the randomized trial arms (\texttt{autoFmtGroup IS NOT NULL}), retaining 556 samples, 4,256 ASVs, and 27,997,124 reads.
  \item \textbf{Longitudinal time window:} Samples were constrained to the window spanning 10 days prior to allo-HCT through 70 days post-HCT (\texttt{DayRelativeToNearestHCT BETWEEN -10 AND 70}), retaining 406 samples, 3,680 ASVs, and 23,069,077 reads.
  \item \textbf{Feature prevalence filtering:} Singleton features appearing in only a single sample were removed (\texttt{--p-min-samples 2}), reducing the feature space to 2,458 ASVs across 406 samples (22,972,681 reads).
  \item \textbf{Taxonomic quality filtering:} Features lacking a phylum-level assignment (\texttt{--p-include p\_\_}) or annotated as plant chloroplast or mitochondrial sequences (\texttt{--p-exclude 'p\_\_;,Chloroplast,Mitochondria'}) were excluded, retaining 406 samples, 2,425 ASVs, and 22,909,139 reads.
  \item \textbf{Sequencing depth threshold:} Samples with fewer than 10,000 total reads (\texttt{--p-min-frequency 10000}) were removed prior to diversity analysis, yielding a final analytical dataset of 356 samples, 2,422 ASVs, and 22,692,526 reads across 24 patients (10 control patients contributing 150 samples and 14 auto-FMT treatment patients contributing 206 samples; 2 to 34 samples per patient, mean 14.8).
\end{enumerate}

\subsection{Upstream sequence denoising and taxonomic classification}
In the upstream demonstration workflow, paired-end sequencing reads from 41 samples were imported in \texttt{CasavaOneEightSingleLanePerSampleDirFmt} format and summarized with \texttt{qiime demux summarize}. Quality filtering, error modeling, paired-end merging, and chimera removal were performed using DADA2 \citep{callahan2016dada2} (\texttt{qiime dada2 denoise-paired}) with forward truncation length set to 204 bp (\texttt{--p-trunc-len-f 204}) and reverse truncation length set to 205 bp (\texttt{--p-trunc-len-r 205}).

Taxonomic classification of representative ASV sequences in the downstream workflow was performed using a multinomial Naive Bayes classifier (\texttt{qiime feature-classifier classify-sklearn}) trained on the SILVA 138.2 99\% small subunit (SSU) rRNA reference database \citep{quast2012silva} (\texttt{silva-138.2-ssu-nr99-515f-806r-classifier.qza}). Reference sequences had been curated via RESCRIPt (\texttt{qiime rescript cull-seqs} with default thresholds of $\ge 5$ degenerate bases and $\ge 8$ bp homopolymers; \texttt{filter-seqs-length-by-taxon} requiring $\ge 900$ bp for Archaea, $\ge 1,200$ bp for Bacteria, and $\ge 1,400$ bp for Eukaryota), extracted for the V4 hypervariable region using primer pair 515F (\texttt{GTGYCAGCMGCCGCGGTAA}) and 806R (\texttt{GGACTACNVGGGTWTCTAAT}), and dereplicated in \texttt{uniq} mode.

\subsection{Phylogenetic reconstruction and alpha/beta diversity analysis}
To support phylogenetic diversity metrics, representative sequences from the 2,422 retained ASVs were aligned with MAFFT, masked for hypervariable positions, and used to reconstruct a midpoint-rooted phylogenetic tree via FastTree (\texttt{qiime phylogeny align-to-tree-mafft-fasttree}). Alpha rarefaction curves for the Shannon index were inspected up to a maximum depth of 33,000 reads (\texttt{qiime diversity alpha-rarefaction --p-max-depth 33000}), and beta rarefaction stability was evaluated using Bray-Curtis dissimilarities with neighbor-joining clustering at 10,000 reads (\texttt{qiime diversity beta-rarefaction --p-sampling-depth 10000}).

Core alpha diversity metrics (observed features/ASVs, Faith's phylogenetic diversity, Shannon entropy, and Pielou's evenness) and beta diversity distance matrices (unweighted UniFrac, weighted UniFrac, Bray-Curtis, and Jaccard) were computed via a single random subsampling without replacement at a uniform depth of 10,000 reads per sample (\texttt{qiime diversity core-metrics-phylogenetic --p-sampling-depth 10000}) \citep{mcmurdie2013waste,raspet2024facilitating}. Principal coordinates analysis (PCoA) was applied to each beta diversity distance matrix to quantify the proportion of variance explained by the leading axes, and unweighted UniFrac distances were additionally embedded via Uniform Manifold Approximation and Projection (\texttt{qiime diversity umap}). Cross-sectional differences in observed ASVs across categorical metadata variables---including time relative to HCT (\texttt{pre}, \texttt{peri}, \texttt{post}), time relative to FMT (\texttt{pre}, \texttt{peri}, \texttt{post}), randomized intervention group (\texttt{control} vs.\ \texttt{treatment}), stool consistency (\texttt{formed}, \texttt{semi-formed}, \texttt{liquid}), and stem cell graft source (\texttt{BM\_unmodified}, \texttt{PBSC\_unmodified}, \texttt{TCD}, \texttt{cord})---were evaluated using non-parametric pairwise Kruskal-Wallis tests with Benjamini-Hochberg false discovery rate (FDR) correction ($q$-values) via \texttt{qiime diversity alpha-group-significance}.

\subsection{Longitudinal linear mixed-effects modeling and random forest regression}
Because repeated fecal sampling from the same patient violates the independence assumption of cross-sectional rank tests \citep{herman2024assessingb,zhou2021linda,zhang2026bayesian}, longitudinal trajectories of microbial richness were modeled using a linear mixed-effects (LME) model (\texttt{qiime longitudinal linear-mixed-effects}) \citep{barrera2025stochastic,fontaine2026locally}. The model was fit via restricted maximum likelihood (REML) with observed ASVs (\texttt{observed\_features}) as the dependent variable, days relative to FMT (\texttt{day-relative-to-fmt}), treatment group (\texttt{autoFmtGroup}), and their interaction (\texttt{day-relative-to-fmt:autoFmtGroup}) as fixed effects, and patient identifier (\texttt{PatientID}) as a random effect across the 356 observations nested within 24 patients.

To identify microbial taxa exhibiting systematic temporal succession across the post-transplant course, ASV counts were collapsed to the genus level (taxonomic level 6; \texttt{qiime taxa collapse --p-level 6}) and filtered conditionally to retain genera with at least 1\% relative abundance in at least 10\% of samples (\texttt{qiime feature-table filter-features-conditionally --p-prevalence 0.1 --p-abundance 0.01}). Supervised random forest regression (\texttt{qiime longitudinal feature-volatility}) was then trained on the filtered genus table to predict the continuous state variable \texttt{week-relative-to-hct} using \texttt{PatientID} as the individual identifier, yielding out-of-bag predictive accuracy metrics (mean squared error, Pearson $r$, $R^2$, and $p$-value) and normalized feature importance scores for the 20 retained genus-level taxa.
